\section{Introduction}
\label{sec:introduction}

Bernoulli bond percolation on an undirected graph $G=(V,E)$ retains
each edge independently with probability $p\in[0,1]$. Retained edges
are called open. The open cluster $C_x$ of a vertex $x$ consists of the
vertices joined to $x$ by finite open paths, including $x$ itself. Write
$x\leftrightarrow A$ when $C_x$ meets a set $A\subseteq V$, and
$x\leftrightarrow y$ when $y\in C_x$. With $\PP_p$ denoting the product
law, the critical probability is
\[
 \pc(G)=\inf\{p\in[0,1]:
       \PP_p(\text{there is an infinite open cluster})>0\}.
\]
A graph is \emph{quasi-transitive} if its automorphism group has finitely
many orbits on vertices. This includes every vertex-transitive graph
and every Cayley graph of a finitely generated group with a finite
generating set.

The criticality problem asks whether an infinite cluster can first
appear at the threshold itself. Benjamini and Schramm formulated the
general quasi-transitive conjecture in their systematic study of
percolation beyond Euclidean lattices
\cite[Conjecture~4]{BenjaminiSchramm1996}.

\begin{theorem}
\label{thm:main}
Let $G$ be an infinite connected locally finite quasi-transitive
undirected graph. If $\pc(G)<1$, then
\[
 \PP_{\pc(G)}(\text{there is an infinite open cluster})=0.
\]
Equivalently, $\PP_{\pc(G)}(|C_x|=\infty)=0$ for every vertex $x$.
\end{theorem}

Theorem~\ref{thm:main} resolves the Benjamini--Schramm criticality
conjecture positively for Bernoulli bond percolation. Its conclusion
concerns all vertex orbits.
The hypothesis $\pc<1$ is necessary: at parameter one every edge is
open, so an infinite connected graph percolates.
In particular, the theorem includes nearest-neighbor bond percolation on
$\Z^d$ for every $d\ge2$; see Corollary~\ref{cor:lattices}.

\subsection{The two subexponential arguments}

For a vertex $x$, let $B(x,r)$ be the graph-distance ball of integer
radius $r$. Put $V_-(r)=\min_x|B(x,r)|$ and
$V_+(r)=\max_x|B(x,r)|$. Quasi-transitivity makes these extrema finite
and compares them up to a bounded shift of $r$. It also implies that
$V_+(r)^{1/r}$ has a limit. A limit greater than one is the exponential
case covered by Hutchcroft. In the other case, the proof divides
according to the following exhaustive alternatives:
\begin{enumerate}[label=(\roman*)]
\item $\log V_-(r)/\log r\to\infty$;
\item $V_-(r)\le r^d$ along an unbounded sequence of radii, for some
fixed finite $d$.
\end{enumerate}
This division requires no classification of intermediate growth.

In case (i), suppose that an infinite cluster occurs at $\pc$.
At a slightly smaller parameter, large finite clusters still occur
with probability bounded away from zero. The volume growth then
places two distinct large clusters within a distance that is small
on the logarithmic scale of their sizes. Repeatedly bisecting a
path between them produces a single adjacent pair. Each bisection
adds at most one cluster and one possible connecting edge. We choose these
edges using the decrease in a search's conditional success
probability, and control the total cost by relative entropy.
Section~\ref{sec:dfs-entropy} establishes the search and entropy
estimates. Sections~\ref{sec:interpolation} and~\ref{sec:superpolynomial} show
that the resulting adjacent pair contradicts the bound in
Proposition~\ref{prop:two-cluster}: across a fixed edge, distinct
clusters each of size at least $s$ occur with probability at most
$Cs^{-1/2}$.

In case (ii), we apply the finitary structure theorem of Tessera and
Tointon \cite[Corollary~2.4]{TesseraTointon2021} to an auxiliary
vertex-transitive graph. Lifting its quotient action to $G$ and
passing to a finite-index subgroup gives a finitely generated
nilpotent quotient and a kernel with finite vertex orbits.
Section~\ref{sec:structure} uses an
induction on the number of infinite cyclic factors in a chosen
series of that group. The one-dimensional case has $\pc=1$.
Otherwise there are two integer coordinates on vertices that change
by a fixed translation under each group element. The induction excludes an infinite open ray whose
two coordinates both remain bounded.

Under the assumed critical percolation, uniqueness and this exclusion
produce two high-probability approximate moves in linearly independent
coordinate directions in Section~\ref{sec:moves}. After rescaling, the moves advance by
$(1,1)$ and $(1,-1)$, with arbitrarily small edge steps.
They persist at a parameter below $\pc$. In
Section~\ref{sec:corridor-exploration}, we place corridors in
these coordinates, with intersections confined to their common
end regions. The joint gluing inequality
(Theorem~\ref{thm:joint-gluing}) transfers connections through each
fresh corridor: reaching a relay set whose vertices have uniformly
high target probabilities gives a high conditional chance of continuing
to the target. Finally an adaptive exploration keeps the
conditional chance of failure small at every tested site. A planar
boundary count gives positive probability of infinitely many
successful sites, hence an infinite open cluster below $\pc$.
All paths and all revealed edges belong to $G$; the quotient supplies
coordinates for this construction.


\subsection{History and related methods}

Percolation was introduced by Broadbent and Hammersley
\cite{BroadbentHammersley1957}. For nearest-neighbor bond percolation
on $\Z^d$, planar methods established critical nonpercolation in
dimension two \cite{Harris1960,Kesten1980}. The lace expansion gave
the result in sufficiently high dimensions \cite{HaraSlade1990},
and Fitzner and van der Hofstad obtained it for $d\ge11$
\cite{FitznerHofstad2017}. The sharpness theorems of Aizenman--Barsky
and Menshikov describe the subcritical regime
\cite{AizenmanBarsky1987,Menshikov1986}; they do not by themselves
settle the behavior at the threshold.

Beyond lattices, Benjamini, Lyons, Peres, and Schramm proved critical
nonpercolation for nonamenable graphs admitting a unimodular
quasi-transitive action \cite{BLPS1999Critical}. Tim\'ar excluded
infinitely many infinite critical clusters on transitive
nonunimodular graphs \cite{Timar2006Nonunimodular}, leaving the
possibility of a unique infinite cluster. Hutchcroft excluded all
infinite critical clusters on every quasi-transitive graph of
exponential growth \cite[Theorem~1.2]{Hutchcroft2016}; we use that
entire theorem. Hermon and Hutchcroft treated unimodular
quasi-transitive graphs satisfying a uniform stretched-exponential
return-probability bound with exponent greater than $1/2$
\cite[Theorem~1.2]{HermonHutchcroft2021}. The two arguments above
cover the remaining subexponential case.

Our joint gluing theorem also has a finite-graph and lattice history
that is distinct from the general quasi-transitive conclusion.
Kozma and Nitzan conjectured the multiplicative connection inequality
\cite[Conjecture~1]{KozmaNitzan2024} and showed that their weaker
uniform near-one gluing Conjecture~3 implies critical nonpercolation
on $\Z^d$ for every $d\ge2$ \cite[Theorem~6]{KozmaNitzan2024}.
Leder's public account reports a Lean formalization of that lattice
conclusion and displays the first-contact comparison (GEN)
\cite[Section~3.4]{Leder2026}. Although the account does not claim
Conjecture~1, (GEN) implies the finite joint inequality: take its
increasing cluster functional to be the indicator of meeting the
target and bound the resulting weighted mean by its minimum.
The subsequent public manuscript \cite{ConnectionInequalities2026}
explicitly records this consequence and stronger finite connection
comparisons. Its public site-percolation extension
\cite{SiteCriticality2026} treats independent labelled hyperedges
and critical site percolation on $\Z^d$ for $d\ge3$. We give a direct
joint-gluing proof with this antecedent credit. These finite and
lattice results do not assert Theorem~\ref{thm:main} for general
quasi-transitive graphs.

The probabilistic tools are developed with their local predecessors
at their points of use. The adjacent-cluster bound adapts the
boundary-discrepancy method of Aizenman, Kesten, and Newman
\cite[Section~3]{AizenmanKestenNewman1987} and Hutchcroft's
ghost-field formulation \cite[Section~3]{Hutchcroft2020Locality}.
Conditional association was developed by van den Berg and Kahn
\cite{VanDenBergKahn2001}; alternating conditional-cluster resampling
is the method of van den Berg, H\"aggstr\"om, and Kahn
\cite{VanDenBergHaggstromKahn2006}. Decision-tree inequalities
\cite{OSSS2005,DuminilCopinRaoufiTassion2019} provide related links
between exploration and pivotality; here the local estimate uses
the martingale of conditional connection probabilities. The entropy
bound accommodates adaptive read sets whose overlap is controlled
only under a comparison law, beyond the fixed-set setting of
Shearer's inequality and its relative-entropy forms
\cite{ChungGrahamFranklShearer1986,MadimanTetali2010}.

The geometric construction adapts the seed and adaptive-block
strategy of Grimmett and Marstrand \cite{GrimmettMarstrand1990}
and the fresh-collar and stored-prediction arguments of Kozma and
Nitzan \cite[Lemma~10 and proof of Theorem~6]{KozmaNitzan2024}.
Critical extinction in each fixed slab
\cite{DuminilCopinSidoraviciusTassion2016Slabs} does not alone give
the full-lattice conclusion. Choosing useful directions by exit
probabilities is related to Martineau and Tassion's parallelogram
construction \cite{MartineauTassion2017}. Here the two-direction
argument and the conditional-failure boundary count are carried out
on the original graph, including when the reveal domains are countable.

The structural reduction runs from Gromov's polynomial-growth
theorem \cite{Gromov1981} through Trofimov's graph and
near-polynomial-growth theorems \cite{Trofimov1985,Trofimov2003}
to Tessera and Tointon's one-scale theorem
\cite[Corollary~2.4]{TesseraTointon2021}. The latter develops the
approximate-group method of Breuillard, Green, and Tao
\cite{BreuillardGreenTao2012} and supplies the quotient action and
stabilizer information used here. Related nilpotent quotients occur
in percolation locality on polynomial-growth graphs
\cite{ContrerasMartineauTassion2023Locality}. In this paper the
quotient supplies coordinates; every connection and exposure uses
original edges.
